\begin{abstract}
Power-system data are commonly classified and released field by field, although their inference risk depends on the other information available to an adversary. A price or renewable forecast that reveals little alone can become informative when combined with another operational series. We develop a background-aware audit that measures this dependence and identifies the combinations responsible for it. Starting from the out-of-sample prediction accuracy of retrained attackers, we define a field's budgeted marginal risk as its largest inferability gain over admissible backgrounds. The resulting profile separates single-field risk from background amplification and records the backgrounds that produce the largest gains. We establish its safety-margin interpretation and connect threshold-crossing fields to minimal unsafe combinations. To evaluate thousands of field sets efficiently, an amortized estimator combines reconstruction-trained and random features with a readout fitted separately for each set. Retraining a few proposed backgrounds then verifies the risk evidence. A controlled dispatch case explains the amplification through unit output, prices and congestion. Across four power-system data sets, ten targets and 124,425 exhaustively retrained field sets, the median average-to-maximum contextual gain ratio is 0.32. Single-field screening recovers only 13--17\% of critical fields; the proposed audit recovers 76--91\% without false positives in the evaluated settings. The estimator takes 48--64\,ms per set and recovers comparable certified marginal risk to TabPFN v2 at substantially lower cost. These results support field-level grading that explicitly accounts for background information.
\end{abstract}
